\section*{Примечания о чтениях}
\subsection*{\texttt{GELFAND-P001475}}
\par\noindent\leavevmode\phantomsection\label{GELFAND-0160-U001}\label{GELFAND-P001475:complex-cancel}\textbf{\texttt{GELFAND-0160-U001}} (complex-cancel).\nopagebreak[4]\par
В скобках перед ∇ad²=0 густо зачёркнут небольшой знак; видны петля и пересекающие штрихи.\par
Кандидаты: \textbackslash{}varphi; \textbackslash{}Phi; неразборчивый зачёркнутый знак. Чтение не установлено в дипломатическом исходнике.\par
Исходные пиксели x=1995–2110, y=885–985; просмотрены оригинал, контекст при 3× и знак при 10×. Исходный знак зачёркиваниями восстановить уверенно нельзя.\par
Повторная сверка: Петля остаётся закрыта зачёркиванием при 5× и 10×; ни один из двух знаков не установлен.\par\medskip
\par\noindent\leavevmode\phantomsection\label{GELFAND-0160-U002}\label{GELFAND-P001475:morse-function-insert}\textbf{\texttt{GELFAND-0160-U002}} (morse-function-insert).\nopagebreak[4]\par
Над «ф-я — морсовская» вставлен интеграл и небольшой v-образный знак.\par
Кандидаты: \textbackslash{}int\textbackslash{}nu; \textbackslash{}int v. Чтение не установлено в дипломатическом исходнике.\par
Исходные пиксели x=580–760, y=1220–1360; просмотрены 4× и 10×. Знак интеграла виден уверенно, греческое ν и латинское v в этой записи не различаются надёжно. Сходный знак виден в (25) на GELFAND-P001477; сравнение не разрешило ν/v. Обозначение не восстановлено по математическому ожиданию.\par
Повторная сверка: Интеграл ясен; следующий знак по-прежнему допускает ν или v. Имя функционала не добавлено.\par\medskip
\par\noindent\leavevmode\phantomsection\label{GELFAND-0160-U003}\label{GELFAND-P001475:homology-equivalence}\textbf{\texttt{GELFAND-0160-U003}} (homology-equivalence).\nopagebreak[4]\par
После (16), где написаны любое представление π₁ и добавление «или ∀ лок. сист.», стоит «Это равносильно тому, что»; в (17) написано M — гомологич. сфера, т.е. H₁(X,ℤ)=0.\par
Кандидаты: Буквально: H\_1(X,\textbackslash{}mathbb Z)=0; Возможная оговорка обозначения: H\_1(M,\textbackslash{}mathbb Z)=0. Чтение не установлено в дипломатическом исходнике.\par
Область (16)–(17): x=160–1790, y=1480–1995; (16) просмотрена 4×, (17) непосредственно в исходном разрешении. Буква X в (17) сохранена. Переход от условия с локальными системами к обычной гомологической сфере требует уточнения гипотез; эта математическая проверка не использована для исправления текста.\par
Повторная сверка: Увеличения подтверждают X и целые коэффициенты. Квантор по всем представлениям не заменён присоединённой локальной системой одной связности; различие рассмотрено отдельно.\par\medskip
\par\noindent\leavevmode\phantomsection\label{GELFAND-0160-U004}\label{GELFAND-P001475:index-theorem-cancel}\textbf{\texttt{GELFAND-0160-U004}} (index-theorem-cancel).\nopagebreak[4]\par
Между «с» и «пом.» в предложении об обычной теореме об индексе густо зачёркнуто слово. Видны лишь части крайних и верхних штрихов.\par
Кандидаты: не записаны. Чтение не установлено в дипломатическом исходнике.\par
Исходные пиксели x=850–1270, y=3040–3210; просмотр при 5×. Ни одна полная буквенная последовательность не читается уверенно; отменённый фрагмент явно отмечен, действующий текст сохранён.\par
Повторная сверка: Две перекрывающиеся детали при 5× и 10× не дают полного чтения слова.\par\medskip
\subsection*{\texttt{GELFAND-P001476}}
\par\noindent\leavevmode\phantomsection\label{GELFAND-0160-U005}\label{GELFAND-P001476:cotangent-cancel}\textbf{\texttt{GELFAND-0160-U005}} (cotangent-cancel).\nopagebreak[4]\par
Между N= и T*L зачёркнут знак, с левой стороны которого видна овальная петля.\par
Кандидаты: \textbackslash{}Omega; неразборчивый зачёркнутый знак. Чтение не установлено в дипломатическом исходнике.\par
Исходные пиксели x=795–940, y=1130–1290; просмотрены 3× и 8×. Не принимать Ω за установленное чтение.\par
Повторная сверка: Овальный остаток совместим с Ω, но зачёркивание не позволяет подтвердить знак.\par\medskip
\par\noindent\leavevmode\phantomsection\label{GELFAND-0160-U006}\label{GELFAND-P001476:below-section-cancel}\textbf{\texttt{GELFAND-0160-U006}} (below-section-cancel).\nopagebreak[4]\par
Под L₁=L зачёркнут короткий математический фрагмент с высоким первым штрихом и низким петлеобразным окончанием.\par
Кандидаты: H\_2; L\_2; неразборчивый зачёркнутый фрагмент. Чтение не установлено в дипломатическом исходнике.\par
Исходные пиксели x=1280–1500, y=1295–1470; просмотрены 3× и 7×. Первоначальные буквы, возможный индекс и зачёркивающие штрихи не отделяются уверенно.\par
Повторная сверка: Верхние и нижние штрихи остаются слиты с зачёркиванием; буква и индекс не восстановлены.\par\medskip
\subsection*{\texttt{GELFAND-P001477}}
\par\noindent\leavevmode\phantomsection\label{GELFAND-0160-U201}\label{GELFAND-P001477:one-form-symbol}\textbf{\texttt{GELFAND-0160-U201}} (one-form-symbol).\nopagebreak[4]\par
В (25), слева в угловых скобках, и в следующей строке перед «морсовская» написан одинаковый v-образный знак с высокой загнутой правой ветвью.\par
Кандидаты: \textbackslash{}nu; v. Чтение не установлено в дипломатическом исходнике.\par
Исходные координаты: (510,120)-(575,205) и (385,440)-(465,540), оба участка просмотрены при 10×; общий контекст (330,25)-(1210,420) при 3×. В теле условно использовано ν; латинское v по начертанию не исключается. На GELFAND-P001475 при повторном просмотре отмечен сходный знак под интегралом. Никакое имя функционала не добавлено; в (25) нет видимого dt.\par
Повторная сверка: Оба знака имеют одинаковую высокую загнутую правую ветвь; ν/v не разрешено. Дифференциал не вставлен.\par\medskip
\par\noindent\leavevmode\phantomsection\label{GELFAND-0160-U202}\label{GELFAND-P001477:decomposition-letter}\textbf{\texttt{GELFAND-0160-U202}} (decomposition-letter).\nopagebreak[4]\par
После теоремы о M и фразы о поверхности Z видна формула W=W₁∪W₂, ∂Wᵢ=Z; первая буква в формуле уверенно читается W.\par
Кандидаты: не записаны. Чтение не установлено в дипломатическом исходнике.\par
Исходный участок (330,1115)-(1430,1290) просмотрен при 3×; контекст теоремы (295,570)-(2150,1130) при 3×. В тексте сохранено W. Неясно, вводится ли иное обозначение или имеется описка; W не заменено на M.\par
Повторная сверка: Первая W уверенно видна при 5× и 10×; связь с предыдущим M не разъяснена в буквальном тексте.\par\medskip
\par\noindent\leavevmode\phantomsection\label{GELFAND-0160-U203}\label{GELFAND-P001477:cancelled-representation-word}\textbf{\texttt{GELFAND-0160-U203}} (cancelled-representation-word).\nopagebreak[4]\par
Между «Предст. π₁(M)» и тире перед «в нем» находится одно слово, перекрытое несколькими жирными горизонтальными штрихами; различим длинный нижний штрих внутри слова.\par
Кандидаты: не записаны. Чтение не установлено в дипломатическом исходнике.\par
Исходный участок (870,1835)-(1090,1970) просмотрен при 6×; строка (300,1800)-(2160,2000) при 3×. Отдельные буквы нельзя надёжно отделить от зачёркивания. В теле сохранён сам факт и место отменённого слова.\par
Повторная сверка: Нижний длинный штрих не устанавливает полного зачёркнутого слова.\par\medskip
\par\noindent\leavevmode\phantomsection\label{GELFAND-0160-U204}\label{GELFAND-P001477:cylinder-decomposition}\textbf{\texttt{GELFAND-0160-U204}} (cylinder-decomposition).\nopagebreak[4]\par
Рисунок (28) показывает две выпуклые крайние части, соединённые двумя почти горизонтальными линиями, и дуги у соединений. Подписи W₁, W₂ и «цилиндр» читаются. У верхних соединений есть короткие косые штрихи.\par
Кандидаты: не записаны. Чтение не установлено в дипломатическом исходнике.\par
Исходный участок (130,1950)-(1150,2270) просмотрен при 3×, затем сопоставлен с полным оригиналом 2408×3480. Геометрическая роль коротких штрихов и слабых дуг не вполне ясна, поэтому рисунок сохранён как SourceFigure без добавления направлений или связей.\par
Повторная сверка: В частичном воссоздании набраны три ясные подписи. Сохранены нижние разрывы; верхние косые штрихи не приняты за уверенные стрелки.\par\medskip
\par\noindent\leavevmode\phantomsection\label{GELFAND-0160-U205}\label{GELFAND-P001477:cancelled-before-path}\textbf{\texttt{GELFAND-0160-U205}} (cancelled-before-path).\nopagebreak[4]\par
Перед ясно написанным «Путь» стоит одно зачёркнутое слово с крупной начальной буквой, похожей на Р; средние и конечные буквы перекрыты двумя горизонтальными линиями.\par
Кандидаты: не записаны. Чтение не установлено в дипломатическом исходнике.\par
Исходный участок (340,2300)-(560,2420) просмотрен при 6×; общий нижний участок (325,2285)-(1980,3100) при 3×. Чтение полного отменённого слова не установлено.\par
Повторная сверка: Р-подобная начальная буква остаётся лишь остатком; незачёркнутое «Путь» передано отдельно.\par\medskip
\par\noindent\leavevmode\phantomsection\label{GELFAND-0160-U206}\label{GELFAND-P001477:four-dimensional-cut}\textbf{\texttt{GELFAND-0160-U206}} (four-dimensional-cut).\nopagebreak[4]\par
В последней скобке ясно написано «M⁴ разрезается на 2 трёхмерных».\par
Кандидаты: не записаны. Чтение не установлено в дипломатическом исходнике.\par
Исходный участок (325,2285)-(1980,3100) просмотрен при 3×; сама скобка находится примерно в (340,2920)-(1700,3070). Размерностное содержание фразы не уточнено в видимом тексте. Все видимые слова и показатель 4 сохранены, без добавления существительного или замены размерности.\par
Повторная сверка: Слова и показатель 4 ясны. Отдельная размерностная проверка не добавляет отсутствующего существительного и не исправляет фразу.\par\medskip
\subsection*{\texttt{GELFAND-P001478}}
\par\noindent\leavevmode\phantomsection\label{GELFAND-0160-U207}\label{GELFAND-P001478:cancelled-margin-label}\textbf{\texttt{GELFAND-0160-U207}} (cancelled-margin-label).\nopagebreak[4]\par
Слева от строки «Тогда есть детерминант» виден замазанный знак пункта и правая круглая скобка.\par
Кандидаты: г); 1). Чтение не установлено в дипломатическом исходнике.\par
Исходный участок (205,1980)-(325,2120) просмотрен при 6×; общий контекст (175,1920)-(2280,2455) при 3×. Нижняя часть знака многократно обведена, поэтому кандидат не принят за чтение.\par
Повторная сверка: Многократная обводка не позволяет выбрать букву или цифру.\par\medskip
\par\noindent\leavevmode\phantomsection\label{GELFAND-0160-U208}\label{GELFAND-P001478:cancelled-final-symbol}\textbf{\texttt{GELFAND-0160-U208}} (cancelled-final-symbol).\nopagebreak[4]\par
В последней строке перед «и т.п.» стоит один густо исправленный или зачёркнутый знак с округлой левой частью и нижним правым хвостом.\par
Кандидаты: a; и. Чтение не установлено в дипломатическом исходнике.\par
Исходный участок (330,3220)-(450,3370) просмотрен при 6×; общий низ (290,2420)-(1690,3360) при 3×. Отдельный знак учтён, но уверенного чтения нет.\par
Повторная сверка: Округлый знак остаётся неясен при 10×.\par\medskip
\subsection*{\texttt{GELFAND-P001479}}
\par\noindent\leavevmode\phantomsection\label{GELFAND-0160-U209}\label{GELFAND-P001479:overwritten-section-label}\textbf{\texttt{GELFAND-0160-U209}} (overwritten-section-label).\nopagebreak[4]\par
Перед «8 букв расположены кубиком» находится знак пункта с двумя тёмными округлыми частями и правой скобкой; есть следы обведения или исправления.\par
Кандидаты: в); г). Чтение не установлено в дипломатическом исходнике.\par
Исходный участок (195,320)-(295,430) просмотрен при 6× и при независимом повторном просмотре при 10×. Оба просмотра предпочитают в), но следы исправления не позволяют исключить г). Нумерация не приведена к ожидаемой последовательности.\par
Повторная сверка: Две округлые части по-прежнему скорее допускают в), но чтение не установлено. Перенумерации нет.\par\medskip
\par\noindent\leavevmode\phantomsection\label{GELFAND-0160-U210}\label{GELFAND-P001479:cancelled-spin-value}\textbf{\texttt{GELFAND-0160-U210}} (cancelled-spin-value).\nopagebreak[4]\par
После σ:= и перед ±1 находится короткая запись, плотно закрытая серией наклонных толстых штрихов.\par
Кандидаты: не записаны. Чтение не установлено в дипломатическом исходнике.\par
Исходный участок (2070,1190)-(2225,1320) просмотрен при 7×; строка и решётка (255,890)-(2375,1590) при 3×. Значение отменённой записи не восстанавливается; читаемое ±1 сохранено отдельно.\par
Повторная сверка: Плотные штрихи не позволяют установить отменённое значение.\par\medskip
\par\noindent\leavevmode\phantomsection\label{GELFAND-0160-U211}\label{GELFAND-P001479:ising-cylinder}\textbf{\texttt{GELFAND-0160-U211}} (ising-cylinder).\nopagebreak[4]\par
Вертикальный цилиндр состоит из верхнего овала, боковых линий и нижней дуги; внутри несколько продольных линий. Внизу есть короткий косой штрих слева и поперечный штрих правее.\par
Кандидаты: не записаны. Чтение не установлено в дипломатическом исходнике.\par
Исходный участок (270,1880)-(1940,2160) просмотрен при 3×; сам рисунок примерно (325,1905)-(535,2150). Соединение внутренних линий с контуром и роль коротких штрихов неясны. SourceFigure сохраняет оригинал как источник рисунка без предположений об ориентации или дополнительных рёбрах.\par
Повторная сверка: Воссозданные контур и внутренние продольные отрезки сохраняют разрывы. Два нижних коротких штриха не превращены в стрелки или новые рёбра.\par\medskip
\par\noindent\leavevmode\phantomsection\label{GELFAND-0160-U212}\label{GELFAND-P001479:cancelled-before-eta}\textbf{\texttt{GELFAND-0160-U212}} (cancelled-before-eta).\nopagebreak[4]\par
Слева от формулы η(τ) находится отменённая запись: начальный угловатый знак напоминает Z, следующие линии перекрыты длинными косыми штрихами.\par
Кандидаты: Z …. Чтение не установлено в дипломатическом исходнике.\par
Исходный участок (355,2230)-(610,2410) просмотрен при 6×; общий формульный блок (340,2180)-(2020,2600) при 3×. Возможное Z не принимается за полное прочтение отменённой записи.\par
Повторная сверка: Z-подобный остаток не устанавливает полного отменённого выражения.\par\medskip
\par\noindent\leavevmode\phantomsection\label{GELFAND-0160-U213}\label{GELFAND-P001479:eta-q-normalization}\textbf{\texttt{GELFAND-0160-U213}} (eta-q-normalization).\nopagebreak[4]\par
На одной строке видны η(τ)=q\textasciicircum{}(1/24)∏(1−qⁿ) и q=e\textasciicircum{}(iπτ); на следующей θ₃(τ)=∑e\textasciicircum{}(iπτn²). Множителя 2 в показателе определения q не видно, пределы произведения и суммы не написаны.\par
Кандидаты: не записаны. Чтение не установлено в дипломатическом исходнике.\par
Исходный формульный блок (340,2180)-(2020,2600) просмотрен при 3×; аргументы τ отдельно в (680,2210)-(895,2530) при 5×. Формулы оставлены буквально. Возможное несогласование нормировок отмечено без вставки 2, пределов, коэффициентов или иных исправлений.\par
Повторная сверка: Первоначальные детали обрезали верх показателя. Исправленные детали при 5× и 10× показывают весь показатель 1/24. Сопоставление нормировок вынесено отдельно.\par\medskip
\par\noindent\leavevmode\phantomsection\label{GELFAND-0160-U214}\label{GELFAND-P001479:last-theta-index}\textbf{\texttt{GELFAND-0160-U214}} (last-theta-index).\nopagebreak[4]\par
В третьем модуле нижней формулы числитель имеет θ с небольшим индексом, наиболее похожим на 1; ранее в строке «и аналог.» ясно записаны θ₂, θ₄.\par
Кандидаты: \textbackslash{}theta\_1; \textbackslash{}theta\_4. Чтение не установлено в дипломатическом исходнике.\par
Исходный участок (1280,2870)-(1405,3005) просмотрен при 10×; независимый повторный просмотр (1235,2860)-(1415,3010) при 10× подтвердил возможность 1. В теле условно оставлено θ₁ с явным маркером. Математическая ожидаемость θ₄ не использована для исправления.\par
Повторная сверка: Нижний индекс всё ещё похож на 1. Исторический вариант 4 сохранён без замены по ожидаемой формуле.\par\medskip
\subsection*{\texttt{GELFAND-P001480}}
\par\noindent\leavevmode\phantomsection\label{GELFAND-0160-U401}\label{GELFAND-P001480:top-torus-cycles}\textbf{\texttt{GELFAND-0160-U401}} (top-torus-cycles).\nopagebreak[4]\par
Тор с внутренним отверстием, дугами α и β; на нижней дуге наконечник влево, справа наконечник вниз-вправо. Часть линий разорвана или сдвоена.\par
Кандидаты: не записаны. Чтение не установлено в дипломатическом исходнике.\par
Оригинальные пиксели: (250,290)-(560,565). Осмотр общего изображения и вырезки 5×. Неуверенно, является ли короткий штрих у верхнего края отверстия самостоятельной стрелкой; пунктирные участки β и контур частично сливаются. SourceFigure сохраняет известные подписи и направления, исходная схема остаётся авторитетной.\par
Повторная сверка: Воссоздание сохраняет две ясные стрелки и подписи. Короткому штриху у верхнего края отверстия направление не приписано; разрывы не заполнены.\par\medskip
\par\noindent\leavevmode\phantomsection\label{GELFAND-0160-U402}\label{GELFAND-P001480:middle-cut-torus}\textbf{\texttt{GELFAND-0160-U402}} (middle-cut-torus).\nopagebreak[4]\par
Разрезанный тор с двумя овальными торцами, β сверху и θ у разреза; у левого внутреннего края плотный угловой штрих.\par
Кандидаты: не записаны. Чтение не установлено в дипломатическом исходнике.\par
Оригинальные пиксели: (275,1365)-(610,1605); общий регион (220,1210)-(2060,1760), увеличение 3×. Контуры разреза видны, но внутренний штрих нельзя уверенно отделить от наконечника стрелки. Не добавлено направление склейки.\par
Повторная сверка: Два овальных торца и ясные β,θ воспроизведены. Левый угловой штрих не получает направления стрелки.\par\medskip
\par\noindent\leavevmode\phantomsection\label{GELFAND-0160-U405}\label{GELFAND-P001480:bottom-point-pair}\textbf{\texttt{GELFAND-0160-U405}} (bottom-point-pair).\nopagebreak[4]\par
Первый индекс пары под последней формулой тесно наложен под z; видны j, горизонтальный штрих и 1, перед j возможна 2. Второй индекс читается 2j.\par
Кандидаты: (z\_\{2j-1\},z\_\{2j\}); (z\_\{j-1\},z\_\{2j\}); (z\_\{2j+1\},z\_\{2j\}). Чтение не установлено в дипломатическом исходнике.\par
Оригинальные пиксели: (260,2970)-(610,3170). Осмотр общего нижнего блока 3× и отдельной вырезки 5×. Индекс не восстановлен по ожидаемой парной нумерации ручек.\par
Повторная сверка: Дополнительные увеличения не разрешили все штрихи первого индекса. Второй индекс 2j ясен.\par\medskip
\par\noindent\leavevmode\phantomsection\label{GELFAND-0160-U406}\label{GELFAND-P001480:bottom-cancelled-word}\textbf{\texttt{GELFAND-0160-U406}} (bottom-cancelled-word).\nopagebreak[4]\par
После «в к-рые» одно слово густо зачёркнуто; далее читается «приклеив. ручки (тонкие)».\par
Кандидаты: не записаны. Чтение не установлено в дипломатическом исходнике.\par
Оригинальные пиксели: (1070,3000)-(1340,3170), контекст (1070,3000)-(1660,3170). Осмотр 3× и 5×. Различимы петли и длинный конечный штрих, но исходное слово под зачёркиванием не восстановлено.\par
Повторная сверка: Оставшиеся петли не дают уверенного полного слова.\par\medskip
\subsection*{\texttt{GELFAND-P001481}}
\par\noindent\leavevmode\phantomsection\label{GELFAND-0160-U407}\label{GELFAND-P001481:top-cft-sets}\textbf{\texttt{GELFAND-0160-U407}} (top-cft-sets).\nopagebreak[4]\par
Овал CFT's, вложенный овал RCFT's и линия от надписи ISING к маленькому кресту или наконечнику внутри вложенного овала.\par
Кандидаты: маленький крестик, отмечающий ISING; наконечник указательной стрелки от ISING. Чтение не установлено в дипломатическом исходнике.\par
Оригинальные пиксели: (330,70)-(840,505). Осмотр всего кадра и вырезки (290,60)-(1350,640), увеличение 3×. Овалы и слова читаемы; окончание указателя около (570,310) состоит из нескольких наложенных штрихов.\par
Повторная сверка: Все три слова набраны. Тёмное окончание указателя оставлено коротким пересекающимся штрихом без утверждения о стрелке или отдельном узле.\par\medskip
\par\noindent\leavevmode\phantomsection\label{GELFAND-0160-U408}\label{GELFAND-P001481:rational-cancelled-prefix}\textbf{\texttt{GELFAND-0160-U408}} (rational-cancelled-prefix).\nopagebreak[4]\par
Перед RATIONAL слева видна I и несколько густо перечёркнутых букв после неё.\par
Кандидаты: I [неразборчивые зачёркнутые буквы]. Чтение не установлено в дипломатическом исходнике.\par
Оригинальные пиксели: (300,500)-(490,610). Осмотр 3× и 7×. Начальная I сохранена отдельно; внутренние буквы перечёркнутого фрагмента не прочитаны и не дополнены по теме.\par
Повторная сверка: Начальная I отделена от неразборчивых букв. Зачёркнутый префикс не дополнен до IRRATIONAL.\par\medskip
\par\noindent\leavevmode\phantomsection\label{GELFAND-0160-U410}\label{GELFAND-P001481:torus-period-convention}\textbf{\texttt{GELFAND-0160-U410}} (torus-period-convention).\nopagebreak[4]\par
В обеих строках о F\_j напечатано от руки τ+2πi; выше θ→θ+2π, внизу страницы T:τ→τ+1.\par
Кандидаты: формулы периода с τ+2πi как написано; T:τ→τ+1 как написано ниже. Чтение не установлено в дипломатическом исходнике.\par
Оригинальные пиксели: (270,770)-(1570,1320) и (950,2950)-(1440,3200). Осмотр 3×. Отмечено возможное изменение соглашения или ошибка источника в периоде τ. Все три видимые записи сохранены без исправления.\par
Повторная сверка: Первая деталь периода была пустой, а первоначальная деталь T показывала строку S. Исправленные детали при 5× и 10× подтверждают и τ+2πi, и T:τ→τ+1. Согласования рассмотрены отдельно.\par\medskip
\par\noindent\leavevmode\phantomsection\label{GELFAND-0160-U411}\label{GELFAND-P001481:monodromy-cancelled-word}\textbf{\texttt{GELFAND-0160-U411}} (monodromy-cancelled-word).\nopagebreak[4]\par
После «M = monodromy matrix.» находится маленькое густо зачёркнутое слово.\par
Кандидаты: не записаны. Чтение не установлено в дипломатическом исходнике.\par
Оригинальные пиксели: (1310,1270)-(1520,1350). Осмотр контекста 3× и вырезки 8×. Начальные петли могут принадлежать нескольким буквам; надёжного полного кандидата нет.\par
Повторная сверка: Петли остаются пересечены зачёркиванием.\par\medskip
\par\noindent\leavevmode\phantomsection\label{GELFAND-0160-U412}\label{GELFAND-P001481:ising-product-conventions}\textbf{\texttt{GELFAND-0160-U412}} (ising-product-conventions).\nopagebreak[4]\par
Источник записывает sqrt(theta₂/eta)=q\textasciicircum{}\{1/8\} Π(1+q\textasciicircum{}n), sqrt(theta₃/eta)=Π(1−q\textasciicircum{}\{n−1/2\}), sqrt(theta₄/eta)=Π(1+q\textasciicircum{}\{n−1/2\}), q=e\textasciicircum{}\{2πiτ\}; нижних границ произведений и дополнительных множителей нет.\par
Кандидаты: не записаны. Чтение не установлено в дипломатическом исходнике.\par
Оригинальные пиксели: (310,1660)-(1600,2200); показатель отдельно (665,1670)-(805,1790). Осмотр 3× и 8×. Показатель 1/8, знаки и отсутствие дополнительных нормировочных множителей воспроизведены буквально. Возможные ошибки или иное соглашение о theta и eta не исправлены; ожидаемые формулы не подставлялись.\par
Повторная сверка: Детали при 5× и 10× подтверждают 1/8, минус в строке θ₃ и плюс в строке θ₄. Стандартные произведения вычислены отдельно с явным соглашением об аргументе и корне.\par\medskip
\subsection*{\texttt{GELFAND-P001482}}
\par\noindent\leavevmode\phantomsection\label{GELFAND-0160-U414}\label{GELFAND-P001482:top-genus-two-twists}\textbf{\texttt{GELFAND-0160-U414}} (top-genus-two-twists).\nopagebreak[4]\par
Двуручная поверхность с подписями α₁,β₁,α₂,β₂,α₃; петли и отрезки пересекают или перекрывают края отверстий.\par
Кандидаты: не записаны. Чтение не установлено в дипломатическом исходнике.\par
Оригинальные пиксели: (330,190)-(900,375), контекст (140,30)-(1050,535), увеличение 3×. Подписи сохранены, но все скрытые/видимые участки кривых и места соединений нельзя восстановить уверенно. Никаких дополнительных пересечений или стрелок не введено.\par
Повторная сверка: Частичное воссоздание сохраняет пять подписей и видимые фрагменты кривых. В неясных наложениях оставлены разрывы; гомологические расчёты не восстанавливают топологию рисунка.\par\medskip
\par\noindent\leavevmode\phantomsection\label{GELFAND-0160-U415}\label{GELFAND-P001482:ln-integration-symbol}\textbf{\texttt{GELFAND-0160-U415}} (ln-integration-symbol).\nopagebreak[4]\par
В двух строках с L\_n одна и та же переменная имеет начертание, похожее на ζ или ξ. Под первым контурным интегралом небольшая c/C. Дифференциал не виден.\par
Кандидаты: переменная ζ, контур c; переменная ξ, контур c; переменная ζ, контур C; переменная ξ, контур C. Чтение не установлено в дипломатическом исходнике.\par
Оригинальные пиксели: (410,2330)-(1850,2680); первый интеграл отдельно (695,2390)-(900,2540). Осмотр 3× и 6×. В основном чтении набрано ζ и c с явным маркером альтернатив; отсутствующий в источнике дифференциал не добавлен.\par
Повторная сверка: Имена ζ/ξ и c/C остаются вариантами. Дифференциал и нормировка вычета не добавлены.\par\medskip
\par\noindent\leavevmode\phantomsection\label{GELFAND-0160-U416}\label{GELFAND-P001482:rho-trailing-fragment}\textbf{\texttt{GELFAND-0160-U416}} (rho-trailing-fragment).\nopagebreak[4]\par
После ρ\_i:A→End H\_i и запятой написано ещё ρ\_i и один округлый знак, похожий на C или O. Продолжения справа не видно.\par
Кандидаты: ρ\_i C; ρ\_i O. Чтение не установлено в дипломатическом исходнике.\par
Оригинальные пиксели: (1100,2740)-(1390,2855), контекст (360,2690)-(1420,2950). Осмотр 3× и 6×. Это видимый короткий исходный фрагмент, а не пропущенное продолжение; его значение не дополнено. Под ним начинается разделительная строка мелких штрихов.\par
Повторная сверка: Короткий фрагмент и разделитель сохранены отдельно; незавершённый смысл не дополнен.\par\medskip
\subsection*{\texttt{GELFAND-P001483}}
\par\noindent\leavevmode\phantomsection\label{GELFAND-0160-U601}\label{GELFAND-P001483:planck-unit}\textbf{\texttt{GELFAND-0160-U601}} (planck-unit).\nopagebreak[4]\par
После слов «планковская длина есть» записано «\textasciitilde{}10\textasciicircum{}\{-33\}м»; перед м отдельной с не видно.\par
Кандидаты: \textasciitilde{}10\textasciicircum{}\{-33\} м. Чтение не установлено в дипломатическом исходнике.\par
Оригинал 2408×3480 px, область (1640,205,1875,285), увеличение 8×; также проверено в полном кадре и текстовой полосе при 3×. Единица м читается уверенно, сохранена. Возможная ошибка источника: указанная степень соответствует распространенной записи длины в сантиметрах, но единица в рукописи выглядит как м; исправление на см не внесено. Неопределенным остается авторский замысел, а не наличие отдельного видимого знака с.\par
Повторная сверка: Видимая единица — м. Отдельное сравнение с планковской длиной в СИ не меняет степень и единицу источника.\par\medskip
\par\noindent\leavevmode\phantomsection\label{GELFAND-0160-U602}\label{GELFAND-P001483:sk-annotation}\textbf{\texttt{GELFAND-0160-U602}} (sk-annotation).\nopagebreak[4]\par
От S\_K проведена короткая черта к трехстрочной помете; первая строка «компактный», вторая похожа на «конд. на», третья оканчивается «пр-ва». Начало третьей строки слито.\par
Кандидаты: конд. на утолщ. пр-ва; конд. на уплотн. пр-ва; конд. на уполщ. пр-ва; конд. на уплощ. пр-ва. Чтение не установлено в дипломатическом исходнике.\par
Оригинальные координаты (1880,955,2280,1230), увеличение 5×, повторно полный кадр. Независимый второй просмотр подтвердил «компактный» и преимущество «конд. на» перед «конф. на», но не разрешил третью строку. Сокращения не раскрыты.\par
Повторная сверка: Первоначальная третья деталь была почти пустой. Исправленные детали при 5× и 10× показывают слитное начало, но не разрешают четыре исторических варианта. Сокращения не раскрыты.\par\medskip
\par\noindent\leavevmode\phantomsection\label{GELFAND-0160-U603}\label{GELFAND-P001483:cft-cancellation}\textbf{\texttt{GELFAND-0160-U603}} (cft-cancellation).\nopagebreak[4]\par
Между «Поэтому» и «это интересно» одно слово зачеркнуто несколькими горизонтальными штрихами; видимые остатки допускают «имеет» или «имеют».\par
Кандидаты: имеет; имеют. Чтение не установлено в дипломатическом исходнике.\par
Оригинальные координаты (510,1940,740,2040), увеличение 7×; полный текст проверен при 3×. Перекрытые буквенные штрихи нельзя отделить от линии отмены; отмена сохранена редакционными квадратными скобками.\par
Повторная сверка: Закрытое окончание остаётся неясным; грамматическое число не подставлено.\par\medskip
\par\noindent\leavevmode\phantomsection\label{GELFAND-0160-U604}\label{GELFAND-P001483:field-overwrite}\textbf{\texttt{GELFAND-0160-U604}} (field-overwrite).\nopagebreak[4]\par
В строке «x\textasciicircum{}β, ... (x) — поле» знак поля перед круглой скобкой густо переписан; видны вертикаль и перекрывающиеся петли.\par
Кандидаты: \textbackslash{}varphi(x); \textbackslash{}Phi(x). Чтение не установлено в дипломатическом исходнике.\par
Оригинальные координаты (1000,2570,1190,2680), увеличение 8×, дополнительно вся нижняя часть при 3×. В body приведено предварительное чтение varphi со связанным маркером Unclear; чтение не утверждается. Верхний индекс β при x отдельно проверен при 8× в (820,2560,940,2670) и прочитан уверенно.\par
Повторная сверка: Переписанный знак поля допускает оба варианта. β читается независимо.\par\medskip
\par\noindent\leavevmode\phantomsection\label{GELFAND-0160-U605}\label{GELFAND-P001483:contour-ending}\textbf{\texttt{GELFAND-0160-U605}} (contour-ending).\nopagebreak[4]\par
После «Φ\{C\} — ф-л от» виден корень «контур» и густо переписанное окончание.\par
Кандидаты: контура; контуров. Чтение не установлено в дипломатическом исходнике.\par
Оригинальные координаты (770,2820,1040,2950), увеличение 6×, общий нижний фрагмент при 3×. Число и падеж окончания не восстанавливались по математическому контексту; буквенный материал окончания учтен в маркере.\par
Повторная сверка: Окончание остаётся неясным; английский текст явно сохраняет альтернативу числа.\par\medskip
\par\noindent\leavevmode\phantomsection\label{GELFAND-0160-U606}\label{GELFAND-P001483:eq1-first-symbol}\textbf{\texttt{GELFAND-0160-U606}} (eq1-first-symbol).\nopagebreak[4]\par
Первый знак поля внутри угловых скобок формулы (1) переписан; остальные однотипные знаки поля читаются как varphi.\par
Кандидаты: \textbackslash{}varphi(x\_1); \textbackslash{}Phi(x\_1). Чтение не установлено в дипломатическом исходнике.\par
Оригинальные координаты (170,2955,740,3105), увеличение 5×, также полный кадр. Первый индекс x\_1 проверен отдельно сравнением с правой частью (1190,2965,1660,3105), 5×. Положительный показатель e\textasciicircum{}\{Σ(φ)\} и отсутствие нормировочного множителя сохранены без дополнений.\par
Повторная сверка: Полные детали показывают x₁ слева и xₙ справа. Третья первоначальная деталь повторяла правый множитель и не служит доказательством x₁; первые две покрывают нужные индексы. Первый знак поля не установлен.\par\medskip
\par\noindent\leavevmode\phantomsection\label{GELFAND-0160-U607}\label{GELFAND-P001483:particle-diagram}\textbf{\texttt{GELFAND-0160-U607}} (particle-diagram).\nopagebreak[4]\par
В нижнем левом углу три линии с подписями x\_1, x\_2, x\_3: верхняя волнистая, две другие поднимаются к ней вправо. В двух местах сближения линий имеются темные штрихи/точки; крайний правый штрих похож также на наконечник.\par
Кандидаты: Две точки соединения трех входящих линий; Одна точка соединения и правый наконечник. Чтение не установлено в дипломатическом исходнике.\par
Оригинальные координаты (180,3105,660,3390), просмотрены как часть (100,2515,2230,3410) при 3× и повторно в полном кадре. Непрерывность линий и характер правого штриха не установлены. Рисунок полностью учтен через SourceFigure, подписи и соседний текст сохранены, топология не дорисована.\par
Повторная сверка: Детали при 10× покрывают оба верхних сближения. Первоначальная нижняя деталь была почти пустой; исправленные детали подтверждают x₃. Три подписи набраны; у неясных контактов оставлены зазоры без утверждения узла или наконечника.\par\medskip
\subsection*{\texttt{GELFAND-P001484}}
\par\noindent\leavevmode\phantomsection\label{GELFAND-0160-U608}\label{GELFAND-P001484:eq2-relation}\textbf{\texttt{GELFAND-0160-U608}} (eq2-relation).\nopagebreak[4]\par
Между правой угловой скобкой и интегралом в (2) густо переписанный знак; отчетливый острый конец обращен вправо, центральные горизонтальные штрихи слиты.\par
Кандидаты: \textbackslash{}longrightarrow; \textbackslash{}Rightarrow; =. Чтение не установлено в дипломатическом исходнике.\par
Оригинальные координаты (827,180,965,325), увеличение 8×, также верхняя текстовая полоса при 3× и полный кадр. Вероятнее стрелка вправо поверх прежнего знака, но тип стрелки и исходный знак не установлены; равенство по аналогии с (1) не подставлено.\par
Повторная сверка: Острый правый конец виден, но слои горизонтальных штрихов не разрешают знак отношения.\par\medskip
\par\noindent\leavevmode\phantomsection\label{GELFAND-0160-U609}\label{GELFAND-P001484:eq5-exponent-overwrite}\textbf{\texttt{GELFAND-0160-U609}} (eq5-exponent-overwrite).\nopagebreak[4]\par
В показателе экспоненты формулы (5) видны минус, переписанная буква, похожая на m, и L; над/через m дополнительные штрихи.\par
Кандидаты: -mL. Чтение не установлено в дипломатическом исходнике.\par
Оригинальные координаты (1080,1260,1215,1340), увеличение 10×; средняя часть при 3×. Окончательный вид напоминает m, но первоначальная буква или характер исправления неразборчивы. В body обозначено переписывание, никаких коэффициентов или множителей не восстановлено.\par
Повторная сверка: Видимый вариант −mL сохранён; предыдущее содержание переписанной буквы не восстановлено.\par\medskip
\par\noindent\leavevmode\phantomsection\label{GELFAND-0160-U610}\label{GELFAND-P001484:eq6-measure}\textbf{\texttt{GELFAND-0160-U610}} (eq6-measure).\nopagebreak[4]\par
После (∂φ)\textasciicircum{}2 первая мера в (6) имеет знак с замкнутой нижней петлей и дугой сверху: ∂\textasciicircum{}\{26\}x. Во втором интеграле явно другое начертание d\textasciicircum{}\{26\}x.\par
Кандидаты: \textbackslash{}partial\textasciicircum{}\{26\}x. Чтение не установлено в дипломатическом исходнике.\par
Оригинальные координаты (997,1945,1183,2063), увеличение 8×, общий фрагмент при 3×. Знак ∂ читается уверенно и оставлен без замены на ожидаемое обозначение дифференциала d. Возможная ошибка обозначения в источнике вынесена сюда; авторский смысл не исправлен.\par
Повторная сверка: Два различных знака меры ясны. Условные вариационные расчёты явно используют обычную меру вместо переопределения рукописного знака.\par\medskip
\par\noindent\leavevmode\phantomsection\label{GELFAND-0160-U611}\label{GELFAND-P001484:eq6-cubic-free}\textbf{\texttt{GELFAND-0160-U611}} (eq6-cubic-free).\nopagebreak[4]\par
В (6) второе слагаемое однозначно +∫φ³d\textasciicircum{}\{26\}x; перед формулой действие названо «опис. своб. частицу».\par
Кандидаты: +\textbackslash{}int\textbackslash{}varphi\textasciicircum{}3d\textasciicircum{}\{26\}x. Чтение не установлено в дипломатическом исходнике.\par
Оригинальные координаты малого символа (1302,1950,1406,2077), увеличение 8×; формула и вводная строка проверены в (100,900,1710,2110) при 3×. Кубическое слагаемое не согласуется с обычным смыслом свободного действия; оно не вычеркнуто в источнике, не удалено и не исправлено. Нерешенным остается возможная математическая ошибка автора.\par
Повторная сверка: Степень 3 и незачёркнутый плюс ясны. Формулировка о свободной частице проверена также в целом оригинале; отдельная деталь формулы не служит доказательством вводной строки. Третья вариация вычислена отдельно.\par\medskip
\par\noindent\leavevmode\phantomsection\label{GELFAND-0160-U612}\label{GELFAND-P001484:vertex-index}\textbf{\texttt{GELFAND-0160-U612}} (vertex-index).\nopagebreak[4]\par
В подписи «V\_N... — вершинные операторы» при N ниже строки стоит маленький петлеобразный знак с высокой правой вертикалью; отдельная точка не различима.\par
Кандидаты: V\_\{N\_j\}; V\_\{N\_d\}. Чтение не установлено в дипломатическом исходнике.\par
Оригинальные координаты (400,2570,584,2710), увеличение 7×, нижняя часть при 3×. Контекст допускает j, но рукописная форма близка к d; предварительное чтение j в body связано с маркером Unclear.\par
Повторная сверка: Новый просмотр при 5× и 10× показывает маленький верхний штрих или точкообразную отметку. Это уточняет историческое наблюдение, не заменяя его; отметка не разрешает j/d.\par\medskip
\par\noindent\leavevmode\phantomsection\label{GELFAND-0160-U613}\label{GELFAND-P001484:branching-diagram}\textbf{\texttt{GELFAND-0160-U613}} (branching-diagram).\nopagebreak[4]\par
Рисунок (8): почти вертикальный левый край и изогнутые верхняя и нижняя границы идут к внутренней почти вертикальной черте. Оттуда вправо видны четыре кривые, две внутренние пересекаются. Слова «ветвления» и рисунок связаны большой стрелкой с наконечником влево.\par
Кандидаты: Пересечение внутренних кривых без дополнительного узла; Соединение внутренних кривых в пересечении. Чтение не установлено в дипломатическом исходнике.\par
Оригинальные координаты фигуры и стрелки (600,2715,1480,2980), просмотр (95,2710,2190,3420) при 3× и полный кадр. Геометрия штрихов описана, но их соединения у внутренней вертикали и значение пересечения не определены. SourceFigure сохраняет рисунок как источник; новых вершин или связей не добавлено.\par
Повторная сверка: Две первоначальные нижние детали показывали главным образом указатель или пустое место. Исправленные детали при 5× и 10× покрывают пересечение и нижние соединения: явный узел не установлен, нижний разрыв у внутренней вертикали сохранён. Большой левый указатель ясен.\par\medskip
\par\noindent\leavevmode\phantomsection\label{GELFAND-0160-U614}\label{GELFAND-P001484:branching-cancellation}\textbf{\texttt{GELFAND-0160-U614}} (branching-cancellation).\nopagebreak[4]\par
После «Теор. струн:» перед «запрещены» короткий фрагмент зачеркнут несколькими штрихами; остается нижняя петля первой буквы и высокий последний штрих.\par
Кандидаты: зап; запр. Чтение не установлено в дипломатическом исходнике.\par
Оригинальные координаты (670,2930,815,3040), увеличение 8×, общая нижняя полоса при 3×. Похоже на оставленное начало следующего слова, но точные отмененные буквы не читаются. Отмена представлена квадратными скобками, без восстановления полного «запрещены».\par
Повторная сверка: Видимые остатки допускают только неуверенные короткие начала, а не полное отменённое слово.\par\medskip
